\documentclass[10pt,a4paper]{amsart}
\usepackage[margin=19mm]{geometry}
\usepackage{amsmath,amssymb,mathtools}
\usepackage[scr=rsfs,cal=euler]{mathalfa}
\usepackage{xfrac}
\usepackage[unicode,hidelinks,bookmarks=false]{hyperref}
\newcommand{\ie}{i.e.\ }
\newcommand{\cf}{cf.\ }
\newcommand{\resp}{respectively\ }
\newcommand{\Subsection}{\subsection}
\providecommand{\nobreakdash}{\nobreak\hbox{-}}
\setlength{\emergencystretch}{2em}
\multlinegap0pt
\usepackage{amssymb}
\usepackage[shortlabels]{enumitem}
\usepackage{mathtools}
\usepackage{bm}
\usepackage[normalem]{ulem}
\usepackage{cancel}
\usepackage{csquotes}
\usepackage{tikz}
\usepackage{tcolorbox}
\newtheorem{lemma}{Lemma}
\newtheorem{prop}{Proposition}
\title{The VC-dimension of strongly regular graphs}
\author{Isabel Byrne, John Byrne, Sebastian M. Cioab\u{a}}
\date{Source: arXiv:2609.10330v1 (CC BY 4.0)}
\begin{document}
\maketitle

\noindent\emph{Source and licence.} The VC-dimension of strongly regular graphs. Version arXiv:2609.10330v1 (CC BY 4.0), distributed by arXiv under the Creative Commons Attribution 4.0 International licence, http://creativecommons.org/licenses/by/4.0/, as recorded in the arXiv licence metadata for this exact version and shown on the abstract page https://arxiv.org/abs/2609.10330v1. Full licence notice: \texttt{licenses/arxiv-2609.10330-CC-BY-4.0.txt}.\par\smallskip

\noindent\textit{Editorial scope.} Counted unit: the article's prop prop (source0.tex lines 893--895). The article's complete original proof of it is delivered with it (source0.tex lines 896--956, one \texttt{proof} environment of 61 lines). No mathematics is written, repaired or re-derived: the statement and the proof are the article's own lines, in the article's own notation, and no retained line is substituted or re-flowed.  Retained uncounted as the source-local context the proof needs: lines 844--855.  Nothing was omitted from any retained span, so the package carries no omission record.  No retained line contains a citation, so the wrapper carries no bibliography and no bibliography entry.  No retained line contains a figure environment, an image-inclusion command or drawing code, so no image is delivered and the package declares no asset dependency.  Editorial formatting only, in addition: an AMS \texttt{amsart} preamble that restates the article's own declarations and macros used by the retained lines, namely the environments \texttt{lemma}, \texttt{prop} on separate counters, together with the standard packages those lines need.  The retained environments print in this document as Lemma 1, Proposition 1 in order of appearance, whereas the article prints the same items with the numbering of its own full document.  The complete native source of the article as distributed in the arXiv e-print for this exact version is bundled unchanged as \texttt{sources/209899/source0.tex}.
\begin{lemma} \label{Lemma Latin square transformations}
Let $L$ be a Latin square and suppose that $L'$ is obtained from $L$ by one of the following operations:
\begin{enumerate}
    \item permuting the rows;
    \item permuting the columns;
    \item permuting the entries;
    \item taking the transpose;
    \item interchanging rows and entries (i.e., putting the entry $i$ in position $(L_{i,j},j)$);
    \item interchanging columns and entries (i.e., putting the entry $j$ in position $(i,L_{i,j})$).
\end{enumerate}
Then $G(L)\simeq G(L')$.
\end{lemma}

\begin{prop}
    The VC-dimension of any Latin square graph is at most 4.
\end{prop}

\begin{proof}
    Suppose to the contrary that $X$ is a shattered set of size 5. We may assume $X\subseteq N((1,1,1))$. Let $R=\{(i,j,\ell)\in X:i=1\}$, $C=\{i,j,\ell\in X:j=1\}$, $E=\{(i,j,\ell\in X:\ell=1\}$ so that $X=R\sqcup C\sqcup E$. Set $r=|R|,c=|C|,e=|E|$. By Lemma \ref{Lemma Latin square transformations}, we may assume $r\ge c\ge e$.

    \textit{Case 1:} $(r,c,e)=(2,2,1)$. We may assume that the shattered set consists of the bold entries in the following principal submatrix
    $$\begin{bmatrix}
        1&\textbf{2}&\textbf{3}\\
        \mathbf{a}&*&*\\
        \mathbf{b}&*&*
    \end{bmatrix}$$
    along with an entry equal to 1, somewhere in the Latin square. Label these vertices $x_2,x_3,x_a,x_b,x_1$. Since $N'(x_2,x_3,x_a,x_b,\overline{x_1})\ne\emptyset$, we must have $\{a,b\}\cap\{2,3\}\ne\emptyset$ and $L_{ij}\in\{a,b\}\cap\{2,3\}$ for some $(i,j)\in\{2,3\}\times\{2,3\}$. Without loss of generality, we have
    $$\begin{bmatrix}
        1&\textbf{2}&\textbf{3}\\
        \mathbf{2}&*&*\\
        \mathbf{b}&*&2
    \end{bmatrix}.$$
    Since $N'(x_2,x_3,x_a,x_1,\overline{x_b})\ne\emptyset$, this forces WLOG
        $$\begin{bmatrix}
        1&\textbf{2}&\textbf{3}&*\\
        \mathbf{2}&3&*&\mathbf{1}\\
        \mathbf{b}&*&2&*
    \end{bmatrix}$$
    and moreover $b\ne 3$, so we may assume $b=4$. Let $(i,j,\ell)\in N'(x_2,x_3,x_1,x_b,\overline{x_a})$. If $i\ne 1$ then we must have $i=2$ and $j=3$, but then $\ell\ne 2$, contradicting $(i,j,\ell)\sim x_2$. Thus, $i=1$, forcing
      $$\begin{bmatrix}
        1&\textbf{2}&\textbf{3}&4\\
        \mathbf{2}&3&*&\mathbf{1}\\
        \mathbf{4}&*&2&*
    \end{bmatrix}.$$
    Now let $(i',j',\ell')\in N'(x_2,x_a,x_b,x_1,\overline{x_3})$. If $\ell'=2$ then we must have $(i',j')=(3,4)$, which is impossible since $L_{3,3}=2$. So $\ell'\ne 2$, so $(i',j',\ell')\sim x_2,x_a$ forces $(i',j')\in\{(1,1),(2,2)\}$, and either case leads to a contradiction.

    \textit{Case 2:} $(r,c,e)=(3,1,1)$. The shattered set takes the form
    $$\begin{bmatrix}
        1&\textbf{2}&\textbf{3}&\textbf{4}\\
        \textbf{a}&*&*&*
    \end{bmatrix}$$
    along with a 1 appearing somewhere in the Latin square. Label the vertices in the shattered set $x_2,x_3,x_4,x_a,x_1$. Note that $N'(x_2,x_3,x_4,x_a,\overline{x_1})\ne\emptyset$ forces 
    $$\begin{bmatrix}
        1&\textbf{2}&\textbf{3}&\textbf{4}&5\\
        \textbf{5}&*&*&*&*
    \end{bmatrix}.$$
    with $x_1$ not in column 5. Since $N'(x_2,x_3,x_4,x_1,\overline{x_a})\ne\emptyset$, we may assume that $L_{1,6}=6$ and $x_1$ belongs to column 6. Then $N'(x_2,x_3,x_a,x_1,\overline{x_4})\ne\emptyset$ forces
    $$\begin{bmatrix}
        1&\textbf{2}&\textbf{3}&\textbf{4}&5&6\\
        \textbf{5}&3&*&*&*&\textbf{1}
    \end{bmatrix}$$
    up to permuting columns $2,3$ and entries $2,3$. Now the facts $N'(x_2,x_4,x_a,x_1,\overline{x_3})\ne\emptyset$ and $N'(x_3,x_4,x_a,x_1,\overline{x_2})\ne\emptyset$ force
    $$\begin{bmatrix}
        1&\textbf{2}&\textbf{3}&\textbf{4}&5&6\\
        \textbf{5}&3&4&2&*&\textbf{1}
    \end{bmatrix}.$$
    But then we see $N(x_2,x_3,x_a,\overline{x_4},\overline{x_1})=\emptyset$, a contradiction.

    \textit{Case 3:} $(r,c,e)=(3,2,0)$. The shattered set takes the form
    $$\begin{bmatrix}
        1&\textbf{2}&\textbf{3}&\textbf{4}\\
        \textbf{a}&*&*&*\\
        \textbf{b}&*&*&*
    \end{bmatrix}.$$
    Since $N'(x_2,x_3,x_4,x_a,\overline{x_b})\ne\emptyset$ and $N'(x_2,x_3,x_4,x_b,\overline{x_a})\ne\emptyset$, we may assume that $a=5$ and $b=6$. Now let $(i,j,\ell)\in N'(x_2,x_3,x_a,x_b,\overline{x_4})$. We have $i\ne 1$, so $\ell\in\{2,3\}$. Also, $j\ne 1$, so $\ell\in\{5,6\}$, a contradiction.

    \textit{Case 4:} $r\ge 4$. Let $x_1,x_2,x_3,x_4$ be 4 elements of the shattered set all in row 1. Then $N(x_1,x_2,x_3,\overline{x_4})=\emptyset$, contradicting assumption.
\end{proof}

\end{document}
